% Pursuing Stacks (cotes 134-2 to 134-8): archive cote 134-8, site page 15 (PDF page 16), typescript page 569.
% typescript. Draft compared: ps7.tex:694-748.
% AI-transcribed by Claude Opus 5.5 (claude-opus-5-5, adaptive thinking, effort high), Message Batches API, batch ps4 (msgbatch_012Aoc1NuskATPsyBeKSYK7J), 2026-11-12; not yet checked by eye.
\scanpage{134-8/15}
\mspage{569}

models in (Cat)).

I am sorry I was a bit confused, when describing $f^!$ as a right
adjoint to $Rf_*$ - I was thinking of the analogy with a map of schemes,
or spaces, which is not only "of finite type" in a suitable sense,
b\xout{y}ut moreover \uline{proper} - in which case, in those duality contexts,
$Rf_*$ is canonically isomorphic with the functor denoted by $f_!$ or $Rf_!$ .
Otherwise, the ("non trivial")"duality theorem" will assert, rather,
that \xout{if} the pair
\[ \over{\unsure{L}}{R}f_! \ , \ f^! \]
is a pair of adjoint functors, just as \inserted{is} the pair
\[ f^* \ , \ Rf_* \quad \text{(or simply } f_* \text{ )} \quad . \]
But even when f is assumed to be proper, the sequence (3) is'nt a
sequence of adjoint functors in the stand\unsure{art|ard} duality contexts, namely
$\hand{Rf}_!$ is by no means left adjoint to $f^*$ \inserted{i.e. $f^*$ \unsure{is'nt} isomorphic to $f^!$} (except in extremely special cases,
\xout{\unsure{such as}} practically I would think only étale maps are OK , which in
the context of (Cat) would correspond to \xout{\unsure{functors isomorphic}} maps in
(Cat) isomorphic to a map $A_{/X} \longrightarrow A$ for X in $A^{\wedge}$, namely \xout{which} maps
which are fibering with discrete fibers). This \uline{does} make an important
discrepancy indeed, between the two kind of contexts - and increases
my perplexity, as to whether or not one should expect a "four variance
duality\xout{"} formalism" to make a sense in (Cat). If so, presumably the
$f_!$ or $Rf_!$ it will involve (perhaps via a suitable notion of"proper"
maps in (Cat), as already \unsure{refered|referred} to earlier (section 70)) will be
different after all from the $Lf_!$ we have been working with lately,
embodying homology properties. \xout{(}But so does $Rf_!$ too , in a rather
strong sense,\xout{(} via the\hand{``}duality theorem\hand{''}!\xout{)}

In the various duality contexts, a basic part is played by the
three particular classes of maps: proper maps, smooth maps, and im-%
mersions , and factorisation of maps into an immersion followed by
either a proper, or a smooth map. In the context of (Cat), there
is a very natural way indeed to define the three class\inserted{es} of maps, as
we'll see in next chapter presumably - so natural indeed, that it is
hard to believe that there may be any other reasonable \over{\illegible{1}}{choice}! One very
striking feature (already mentioned in section 70) is that the two \inserted{first}
notions are "dual" to each other in the rather tautological sense,
namely \inserted{that} a map $f: A \longrightarrow B$ in (Cat) is proper iff the corresponding map
$f^o : A^o \longrightarrow B^o$ is smooth (whereas the notion of an immersion is auto-%

\margin{15}
